\begin{siglas}
  \item[ANP] Agência Nacional do Petróleo, Gás Natural e Biocombustíveis
  \item[ASCII] \textit{American Standard Code for Information Interchange} — código padrão americano para intercâmbio de informações
  \item[CENPES] Centro de Pesquisas e Desenvolvimento Leopoldo Américo Miguez de Mello
  \item[CSV] \textit{Comma-Separated Values} — valores separados por vírgulas
  \item[DLIS] \textit{Digital Log Interchange Standard} — padrão de intercâmbio de dados digitais de perfilagem
  \item[EPSG] \textit{European Petroleum Survey Group} — identificação do registro de sistemas de referência de coordenadas
  \item[EW] \textit{East--West} — leste--oeste
  \item[H/V] Razão horizontal/vertical
  \item[HVSR] \textit{Horizontal-to-Vertical Spectral Ratio} — razão espectral horizontal/vertical
  \item[IBGE] Instituto Brasileiro de Geografia e Estatística
  \item[IDW] \textit{Inverse Distance Weighting} — ponderação pelo inverso da distância
  \item[IL] \textit{Inline} — linha do levantamento sísmico
  \item[JCGM] \textit{Joint Committee for Guides in Metrology} — Comitê Conjunto para Guias em Metrologia
  \item[LOC] Laboratório de Oceanografia Costeira
  \item[MTC] \textit{Mass-Transport Complex} — complexo de transporte de massa
  \item[MTD] \textit{Mass-Transport Deposit} — depósito de transporte de massa
  \item[NOAA] \textit{National Oceanic and Atmospheric Administration} — Administração Nacional Oceânica e Atmosférica
  \item[NS] \textit{North--South} — norte--sul
  \item[OBS] \textit{Ocean-Bottom Seismometer} — sismógrafo de fundo oceânico
  \item[ON] Observatório Nacional
  \item[pBC] \textit{Paleo-Brazil Current} — paleocorrente do Brasil
  \item[PCR-BS] Projeto de Caracterização Regional da Bacia de Santos
  \item[PEG-BR] \textit{Pool} de Equipamentos Geofísicos do Brasil
  \item[PPSD] \textit{Probabilistic Power Spectral Density} — densidade espectral de potência probabilística
  \item[PSTM] \textit{Prestack Time Migration} — migração em tempo pré-empilhamento
  \item[RMS] \textit{Root Mean Square} — raiz da média dos quadrados
  \item[RMSE] \textit{Root Mean Square Error} — raiz do erro quadrático médio
  \item[SAD69] \textit{South American Datum 1969} — datum sul-americano de 1969
  \item[SEE] Sucessão Evaporítica Estratificada
  \item[SEG-Y] Formato de intercâmbio de dados sísmicos da \textit{Society of Exploration Geophysicists}
  \item[SEI] Sistema Eletrônico de Informações
  \item[SESAME] \textit{Site Effects Assessment Using Ambient Excitations} — avaliação de efeitos de sítio por meio de excitações ambientais
  \item[SIRGAS] Sistema de Referência Geocêntrico para as Américas
  \item[SOC] \textit{Southern Ocean Current} — corrente do Oceano Austral
  \item[SRC] Sistema de Referência de Coordenadas
  \item[TWT] \textit{Two-Way Traveltime} — tempo duplo de percurso do sinal sísmico
  \item[UFRGS] Universidade Federal do Rio Grande do Sul
  \item[UFSC] Universidade Federal de Santa Catarina
  \item[UTM] \textit{Universal Transverse Mercator} — projeção Universal Transversa de Mercator
  \item[VSP] \textit{Vertical Seismic Profile} — perfil sísmico vertical
  \item[XL] \textit{Crossline} — linha transversal do levantamento sísmico
\end{siglas}
\cleardoublepage
